\qhead{1. Your company — Describe your company, the main products sold, the type of clients (B2B, B2C, local, worldwide, etc.) and the main competitors. Indicate the product(s) the answers focus on.}

BYD (``Build Your Dreams'') began as a battery maker in 1995 and is now the world's largest producer of new-energy vehicles, overtaking Tesla in battery-electric volume in 2025 with about 2.25 million units and revenue near RMB~804\,bn \parencite{byd2025}. Its defining asset is \textbf{extreme vertical integration}: roughly 70--80\% of each car is made in-house, giving an estimated \textbf{15\% cost advantage over Tesla} \parencite{motleyfool2025}. This cost engine underpins almost every pricing decision below.

\textbf{Clients.} BYD is overwhelmingly \textbf{B2C} and \textbf{worldwide} — private buyers in China (its core market) and increasingly Europe, sold through dealers and its own European showrooms.

\textbf{Product ladder.} It spans a Good--Better--Best range from the sub-RMB-60{,}000 Seagull, through the mass-market Dynasty and Ocean series, up to the luxury Yangwang brand (the U9 reaching about \$3\,million) \parencite{electrekU9}.

\textbf{Focus product.} The answers below focus on the \textbf{BYD Qin Plus DM-i}, the mass-market plug-in hybrid marketed as ``electricity cheaper than oil'' \parencite{carscoopsQin}, using the wider ladder where segmentation and price-level questions require it.

\textbf{Main competitors.} Tesla (its direct global rival), the legacy petrol makers it is converting (Toyota, Volkswagen), and the 100-plus domestic Chinese NEV brands locked in an intense price war.
